%
% listener_locks.tex
%
% Companion to docs/listener_lifecycle.tex: the LOCK ORDER the fixed design
% requires, checked for deadlock-freedom.
%
% listener_lifecycle.tex makes each critical section one atomic step, so its
% deadlocked:0 proves nothing about locks (Want and LockExpire are enabled in
% every state).  The claim that the fixed design cannot deadlock is made here,
% and here alone.  The locks in play (listener_lifecycle.tex, Decision 5 and
% the section "Deadlock-freedom, and the lock order the model assumes"):
%
%   rp, rs, rq  the request locks, one per reply kind: pong (ping), ack
%               (show), query.  A request holds its kind's lock from the
%               write to its consumption or failure.
%   st          the generation state lock: curGen, held, started, detaching a
%               sink.  Field updates only; never I/O, never a dial.
%   lk          the connection's _lock: one sendall or one recv.
%   qu          a reply-queue / ack-ledger lock: one put, get or wake-up.
%               The ledger entry is appended under lk, in the write's own
%               critical section.  The queues are modelled as one lock: two
%               queue locks are never held together, so one stands for all.
%
% THE FIXED ORDER:  request lock < st < lk < qu.  The edges that occur are
% request -> lk, lk -> qu, request -> st, request -> qu, st -> qu.  A rank
% function over the locks is strictly increasing along every edge, so the
% acquisition graph is acyclic.
%
% The model is not that argument.  It runs the six programs below, each as a
% thread of acquire / release / I/O steps over the locks, against every
% interleaving, and asks ProB for a state in which some thread is not done and
% none can move.
%
% MUTUAL EXCLUSION IS ASSUMED, AS A LOCK IS MODELLED: Acquire is guarded by
% "the lock is free".  That is what a lock is; it is not checked.  What is
% checked:
%   (1) deadlock-freedom: Halt is enabled only when every thread is done, so a
%       blocked subset is reported by probcli as a deadlock;
%   (2) ORDER, as a negated goal: the flag orderBroken is set by an Acquire
%       made while the thread holds a lock of equal or higher rank.  The goal
%       orderBroken = fset must not be found.  The order is deliberately not a
%       conjunct of the state schema: ProB treats the primed state schema as
%       part of each operation's guard, so a violation written there disables
%       the operation instead of reporting it.
%
% THE WEDGE.  A display that stops reading holds a sender inside sendall,
% under lk, until the connection is shut down.  Wedge models that (it is
% enabled only while a reporter is running, so a wedge always has someone who
% can end it); Write, the I/O under lk, is blocked until shut = fset.  The
% reporter must reach its shutdown without waiting for anything a wedged
% sender holds: it takes st and qu, releases both, and shuts the socket
% outside the state lock.  A reporter that needed lk or a request lock would
% wait on the sender that waits on the shutdown; the model has no such
% reporter, and the deadlock check shows the one it has is never stuck.
%
% WHAT ENDS A WEDGE IN THE CODE.  Not a reporter: a display that stops reading
% produces no end-of-file, so the reader never reports.  The bound is the send
% timeout: every generation's socket carries SO_SNDTIMEO (send_timeout.py,
% one second, applied as one to two seconds by platform, on both legs), a timed-out sendall raises, and the write-failure
% rule (listener_lifecycle.tex, Decision 10(a)) retires the generation.  The
% model's restriction of Wedge to a running reporter stands for that bound.
% Without the bound the model deadlocks: with the reporter condition dropped,
% ProB finds  Start(showC,T2); Wedge; Start(queryC,T1); Acquire(T2);
% Acquire(T1); Acquire(T2)  (the show wedged inside its write holding lk, the
% query waiting on lk for ever).  That is a deadlock of the environment, not
% of the order; it shows why the send timeout must be on every socket.
%
% LOCKS OUTSIDE THE MODEL.  ClientRegistry._lock and the store lock sit above
% the model's set (registry lock -> lk, registry lock -> st).  No lock of this
% model is held while taking either of them; in particular a query's timeout
% retire is close_if_current(g) inside the link and never calls
% ClientRegistry.drop() while holding rq.
%
% THE PROGRAMS (the number of steps is the index of the program counter):
%   pingC   rp; lk{write}; qu{wait}; release rp
%   showC   rs; lk{write; qu{ledger append}}; qu{wait}; release rs
%   queryC  rq; lk{write}; qu{wait}; st{qu{retire the generation, wake
%           waiters}}; shutdown outside st; release rq
%   sendS   the replicator's send_scene: lk{write; qu{ledger append}}
%   reader  the listener's poll_chunk, and equally the inline read of rule 23
%           (select outside lk, recv under it: lk alone, for the recv only);
%           then qu{route the reply}
%   report  st{qu{retire}}; shutdown outside st
% A wait releases its queue lock inside Condition.wait; it is the acquire /
% release pair on qu.  The marker-like leaves that hold nothing else are not
% modelled.
%
% THE REJECTED ORDER (listener_locks_inverted_buggy.tex): showC takes the
% ledger/queue lock before lk.  A reader holds lk and wants qu; showC holds qu
% and wants lk.  ProB finds the deadlock.
%
% Type-check:  fuzz docs/listener_locks.tex
% Model-check: probcli docs/listener_locks.tex -model_check -coverage \
%                -p MAX_OPERATIONS 200 -card T 2
%              probcli docs/listener_locks.tex -model_check -coverage \
%                -p MAX_OPERATIONS 200 -card T 3
%              probcli docs/listener_locks.tex -model_check -nodead -bf \
%                -goal "orderBroken = fset" -p MAX_OPERATIONS 200 -card T 3
% Expected: no counter example, deadlocked 0, goal not found.
%

\documentclass[a4paper,10pt,fleqn]{article}
\usepackage[margin=1in]{geometry}
\usepackage{fuzz}

\begin{document}

\title{lux Listener Lifecycle: the Lock Order}
\maketitle

\begin{zed}
  [T]
\end{zed}

A thread: a caller (the liveness ping, the replicator's show, a query), the
replicator's \texttt{send\_scene}, the reader, the report.  Two suffice for a
cycle; three give every pair a third party.

\begin{zed}
  LOCK ::= rp | rs | rq | lk | st | qu
\end{zed}

\begin{zed}
  ACT ::= grab | drop | io | sd | fin
\end{zed}

\begin{zed}
  PROG ::= pingC | showC | queryC | sendS | reader | report
\end{zed}

\begin{zed}
  FSTATE ::= fclear | fset
\end{zed}

\begin{zed}
  TST ::= tnew | trun | tdone
\end{zed}

Steps: acquire or release a lock, do I/O under the connection's lock, shut the
socket down, finish.

\begin{axdef}
  lvl : LOCK \fun 1 \upto 4
\where
  lvl = \{ rp \mapsto 1, rs \mapsto 1, rq \mapsto 1, st \mapsto 2,
    lk \mapsto 3, qu \mapsto 4 \}
\end{axdef}

The order of the fixed design: a request lock, then the state lock, then the
connection's \texttt{\_lock}, then a queue lock.  The order is a ranking, not
a premise of the model: it is the standard the flag \texttt{orderBroken}
measures each acquisition against.

\begin{axdef}
  act : PROG \cross (1 \upto 13) \pfun ACT \\
  lck : PROG \cross (1 \upto 13) \pfun LOCK
\where
  act = \{ (pingC, 1) \mapsto grab, (pingC, 2) \mapsto grab, (pingC, 3) \mapsto io,
    (pingC, 4) \mapsto drop, (pingC, 5) \mapsto grab, (pingC, 6) \mapsto drop,
    (pingC, 7) \mapsto drop, (pingC, 8) \mapsto fin,
    (showC, 1) \mapsto grab, (showC, 2) \mapsto grab, (showC, 3) \mapsto io,
    (showC, 4) \mapsto grab, (showC, 5) \mapsto drop, (showC, 6) \mapsto drop,
    (showC, 7) \mapsto grab, (showC, 8) \mapsto drop, (showC, 9) \mapsto drop,
    (showC, 10) \mapsto fin,
    (queryC, 1) \mapsto grab, (queryC, 2) \mapsto grab, (queryC, 3) \mapsto io,
    (queryC, 4) \mapsto drop, (queryC, 5) \mapsto grab, (queryC, 6) \mapsto drop,
    (queryC, 7) \mapsto grab, (queryC, 8) \mapsto grab, (queryC, 9) \mapsto drop,
    (queryC, 10) \mapsto drop, (queryC, 11) \mapsto sd, (queryC, 12) \mapsto drop,
    (queryC, 13) \mapsto fin,
    (sendS, 1) \mapsto grab, (sendS, 2) \mapsto io, (sendS, 3) \mapsto grab,
    (sendS, 4) \mapsto drop, (sendS, 5) \mapsto drop, (sendS, 6) \mapsto fin,
    (reader, 1) \mapsto grab, (reader, 2) \mapsto drop, (reader, 3) \mapsto grab,
    (reader, 4) \mapsto drop, (reader, 5) \mapsto fin,
    (report, 1) \mapsto grab, (report, 2) \mapsto grab, (report, 3) \mapsto drop,
    (report, 4) \mapsto drop, (report, 5) \mapsto sd, (report, 6) \mapsto fin \} \\
  lck = \{ (pingC, 1) \mapsto rp, (pingC, 2) \mapsto lk, (pingC, 4) \mapsto lk,
    (pingC, 5) \mapsto qu, (pingC, 6) \mapsto qu, (pingC, 7) \mapsto rp,
    (showC, 1) \mapsto rs, (showC, 2) \mapsto lk, (showC, 4) \mapsto qu,
    (showC, 5) \mapsto qu, (showC, 6) \mapsto lk, (showC, 7) \mapsto qu,
    (showC, 8) \mapsto qu, (showC, 9) \mapsto rs,
    (queryC, 1) \mapsto rq, (queryC, 2) \mapsto lk, (queryC, 4) \mapsto lk,
    (queryC, 5) \mapsto qu, (queryC, 6) \mapsto qu, (queryC, 7) \mapsto st,
    (queryC, 8) \mapsto qu, (queryC, 9) \mapsto qu, (queryC, 10) \mapsto st,
    (queryC, 12) \mapsto rq,
    (sendS, 1) \mapsto lk, (sendS, 3) \mapsto qu, (sendS, 4) \mapsto qu,
    (sendS, 5) \mapsto lk,
    (reader, 1) \mapsto lk, (reader, 2) \mapsto lk, (reader, 3) \mapsto qu,
    (reader, 4) \mapsto qu,
    (report, 1) \mapsto st, (report, 2) \mapsto qu, (report, 3) \mapsto qu,
    (report, 4) \mapsto st \}
\end{axdef}

The programs, one row of \texttt{act} and \texttt{lck} per step.  A ping holds
its request lock across a write under \texttt{\_lock} and a wait on the queue;
a show appends its ledger entry (queue lock) inside the same \texttt{\_lock}
hold as the write, so the ledger's order is the wire's; a query that times out
retires its generation, taking the state lock and, inside it, the queue lock
to wake the sink's waiters, and closes the socket after releasing the state
lock; the reader's \texttt{recv} holds \texttt{\_lock} alone and routes the
reply under a queue lock afterwards; the report retires under the state lock
and shuts the socket down outside it.

\begin{schema}{LS}
  ts : T \fun TST \\
  pr : T \fun PROG \\
  pc : T \fun 1 \upto 13 \\
  owner : LOCK \pfun T \\
  wedged : FSTATE \\
  shut : FSTATE \\
  orderBroken : FSTATE
\end{schema}

\texttt{owner} is the lock table; a lock in its domain is held.  Nothing else
is written into the state schema: no property is a conjunct of it.

\begin{schema}{Init}
  LS'
\where
  ts' = T \cross \{ tnew \} \\
  pr' = T \cross \{ reader \} \\
  pc' = T \cross \{ 1 \} \\
  owner' = \emptyset \\
  wedged' = fclear \\
  shut' = fclear \\
  orderBroken' = fclear
\end{schema}

\begin{schema}{Start}
  \Delta LS \\
  t? : T \\
  p? : PROG
\where
  ts~t? = tnew \\
  ts' = ts \oplus \{ t? \mapsto trun \} \\
  pr' = pr \oplus \{ t? \mapsto p? \} \\
  pc' = pc \\
  owner' = owner \\
  wedged' = wedged \\
  shut' = shut \\
  orderBroken' = orderBroken
\end{schema}

A thread picks its program once.

\begin{schema}{Wedge}
  \Delta LS
\where
  wedged = fclear \\
  shut = fclear \\
  \exists t : T | ts~t = trun \land pr~t = report \spot true \\
  wedged' = fset \\
  ts' = ts \\
  pr' = pr \\
  pc' = pc \\
  owner' = owner \\
  shut' = shut \\
  orderBroken' = orderBroken
\end{schema}

The display stops reading.  The wedge is possible only while a reporter is
running and has not yet shut the socket down: a wedge with nobody to end it
would be a deadlock of the environment, not of the locks.

What ends a wedge in the code is not a reporter.  A display that stops
reading produces no end-of-file, so the reader never reports.  The bound is the
send timeout: every generation's socket carries \texttt{SO\_SNDTIMEO} (one second,
applied as one to two seconds by platform, on both legs), a timed-out \texttt{sendall} raises, and the
write-failure rule retires the generation.  The restriction of \texttt{Wedge}
to a running reporter stands for that bound.  Drop the restriction and ProB
finds the deadlock \texttt{Start(showC,T2)}, \texttt{Wedge},
\texttt{Start(queryC,T1)}, \texttt{Acquire(T2)}, \texttt{Acquire(T1)},
\texttt{Acquire(T2)}: the show wedged inside its write holding
\texttt{\_lock}, the query waiting on \texttt{\_lock} for ever.  It is a
deadlock of the environment, not of the order, and it is why the bound is
needed.

Two locks lie above the model's set and outside it: the
\texttt{ClientRegistry} lock and the store lock (registry lock $\to$
\texttt{\_lock}, registry lock $\to$ state lock).  No lock of this model is held
while taking either of them; in particular a query's timeout retire is
\texttt{close\_if\_current(g)} inside the link, and never calls
\texttt{ClientRegistry.drop()} while holding the request lock.

\begin{schema}{AcquireOrdered}
  \Delta LS \\
  t? : T
\where
  ts~t? = trun \\
  act~(pr~t?, pc~t?) = grab \\
  lck~(pr~t?, pc~t?) \notin \dom owner \\
  \forall l : \dom owner | owner~l = t? \spot
    lvl~l < lvl~(lck~(pr~t?, pc~t?)) \\
  owner' = owner \cup \{ lck~(pr~t?, pc~t?) \mapsto t? \} \\
  pc' = pc \oplus \{ t? \mapsto pc~t? + 1 \} \\
  ts' = ts \\
  pr' = pr \\
  wedged' = wedged \\
  shut' = shut \\
  orderBroken' = orderBroken
\end{schema}

\begin{schema}{AcquireUnordered}
  \Delta LS \\
  t? : T
\where
  ts~t? = trun \\
  act~(pr~t?, pc~t?) = grab \\
  lck~(pr~t?, pc~t?) \notin \dom owner \\
  \exists l : \dom owner | owner~l = t? \spot
    lvl~l \geq lvl~(lck~(pr~t?, pc~t?)) \\
  owner' = owner \cup \{ lck~(pr~t?, pc~t?) \mapsto t? \} \\
  pc' = pc \oplus \{ t? \mapsto pc~t? + 1 \} \\
  ts' = ts \\
  pr' = pr \\
  wedged' = wedged \\
  shut' = shut \\
  orderBroken' = fset
\end{schema}

\begin{zed}
  Acquire \defs AcquireOrdered \lor AcquireUnordered
\end{zed}

Acquire is guarded by ``the lock is free'': that is mutual exclusion, assumed
as a lock is modelled and not checked.  The two disjuncts differ only in the
flag: an acquisition made while the thread holds a lock of equal or higher
rank sets \texttt{orderBroken}, and both proceed, so that a broken order is
reported by the goal instead of hidden by a disabled step.

\begin{schema}{Release}
  \Delta LS \\
  t? : T
\where
  ts~t? = trun \\
  act~(pr~t?, pc~t?) = drop \\
  owner' = \{ lck~(pr~t?, pc~t?) \} \ndres owner \\
  pc' = pc \oplus \{ t? \mapsto pc~t? + 1 \} \\
  ts' = ts \\
  pr' = pr \\
  wedged' = wedged \\
  shut' = shut \\
  orderBroken' = orderBroken
\end{schema}

\begin{schema}{Write}
  \Delta LS \\
  t? : T
\where
  ts~t? = trun \\
  act~(pr~t?, pc~t?) = io \\
  wedged = fclear \lor shut = fset \\
  pc' = pc \oplus \{ t? \mapsto pc~t? + 1 \} \\
  owner' = owner \\
  ts' = ts \\
  pr' = pr \\
  wedged' = wedged \\
  shut' = shut \\
  orderBroken' = orderBroken
\end{schema}

The write under \texttt{\_lock} completes unless the display is wedged, and
a wedged one completes it only once the connection is shut down.  The thread
that waits here holds whatever it holds.

\begin{schema}{Shutdown}
  \Delta LS \\
  t? : T
\where
  ts~t? = trun \\
  act~(pr~t?, pc~t?) = sd \\
  shut' = fset \\
  pc' = pc \oplus \{ t? \mapsto pc~t? + 1 \} \\
  owner' = owner \\
  ts' = ts \\
  pr' = pr \\
  wedged' = wedged \\
  orderBroken' = orderBroken
\end{schema}

The socket's \texttt{shutdown} and \texttt{close} take no lock of the model:
they happen outside the state lock and unblock a wedged \texttt{sendall}.

\begin{schema}{Finish}
  \Delta LS \\
  t? : T
\where
  ts~t? = trun \\
  act~(pr~t?, pc~t?) = fin \\
  ts' = ts \oplus \{ t? \mapsto tdone \} \\
  pc' = pc \\
  owner' = owner \\
  pr' = pr \\
  wedged' = wedged \\
  shut' = shut \\
  orderBroken' = orderBroken
\end{schema}

\begin{schema}{Halt}
  \Xi LS
\where
  \forall t : T \spot ts~t = tdone
\end{schema}

Halt is enabled only when every thread has finished, so a state in which some
thread is not done and no thread can move is a deadlock of the model, and
probcli reports it.  Deadlock-freedom is the claim that every reachable state
other than the finished one has an enabled step: a thread holding a queue
lock waits for nothing; a thread holding \texttt{\_lock} waits only for a
queue lock or, in the write, for a shutdown that the reporter reaches holding
nothing a wedged sender holds; a thread holding the state lock waits only for
a queue lock; a thread holding a request lock waits for the locks above it.
The acquisition graph is acyclic, and the wedge adds no edge into the report.

\begin{verbatim}
  fuzz docs/listener_locks.tex
  probcli docs/listener_locks.tex -model_check -coverage \
    -p MAX_OPERATIONS 200 -card T 2
  probcli docs/listener_locks.tex -model_check -coverage \
    -p MAX_OPERATIONS 200 -card T 3
  probcli docs/listener_locks.tex -model_check -nodead -bf \
    -goal "orderBroken = fset" -p MAX_OPERATIONS 200 -card T 3
\end{verbatim}

At two threads: 3267 states, 6433 transitions, no counter example,
\texttt{deadlocked:0}, ten of ten operations covered.  At three: 150552
states, 417164 transitions, no counter example, \texttt{deadlocked:0}, all
covered.  The goal \texttt{orderBroken = fset}, searched at three threads with
the deadlock check off, is not found in 150551 states analysed.  The control,
\texttt{listener\_locks\_inverted\_buggy.tex}, deadlocks and reaches the goal.

\end{document}
